\section{Introduzione}
Il presente documento rappresenta la relazione individuale relativa al progetto \textit{JTrash}, realizzato nell'ambito del corso \textbf{Metodologie di Programmazione}. L'obiettivo di questo progetto è di applicare le competenze acquisite durante il corso nell'implementazione di un gioco di carte, conformemente alle specifiche fornite.

\subsection{Informazioni Personali}
%\addcontentsline{toc}{subsection}{1.1 Informazioni Personali}
Il sottoscritto, \textbf{Lorenzo Sabatino} iscritto al corso di \textbf{Metodologie di Programmazione} in modalità di studio \textbf{presenza MZ} con numero di matricola \textsf{\colorbox{black}{\textcolor{black}{0000000}}}, ha sviluppato il progetto \textit{JTrash} rispettando le linee guida e le specifiche indicate. Questa relazione dettaglierà l'implementazione, le decisioni di progettazione, l'uso dei design pattern e altre scelte effettuate per la creazione del gioco.

\subsection{Obiettivo del Progetto}
%\addcontentsline{toc}{subsection}{1.2 Obiettivo del Progetto}
L'obiettivo principale di \textit{JTrash} consiste nello sviluppo di un gioco di carte interattivo, che permetta ai giocatori di sfidarsi in partite contro avversari artificiali. Il gioco include funzionalità di gestione del profilo utente, quali la personalizzazione del nickname e dell'avatar, oltre al tracciamento delle partite giocate e dei risultati ottenuti. L'aspirazione è quella di creare un'esperienza coinvolgente e divertente, mantenendo contemporaneamente un'architettura del software e un codice puliti.
